\documentclass[11pt, a4paper]{article}

% --- UNIVERSAL PREAMBLE BLOCK ---
\usepackage[a4paper, top=2.5cm, bottom=2.5cm, left=2cm, right=2cm]{geometry}
\usepackage{fontspec}

\usepackage[chinese, bidi=basic, provide=*]{babel}

\babelprovide[import, onchar=ids fonts]{chinese}
\babelprovide[import, onchar=ids fonts]{english}

% Set default/Latin font to Sans Serif in the main (rm) slot
\babelfont{rm}{Noto Sans}
% Assign a specific font for Chinese text
\babelfont[chinese]{rm}{Noto Sans CJK SC}

% Add because main language is not English
\usepackage{enumitem}
\setlist[itemize]{label=-}

% --- CHINESE TYPOGRAPHY PACKAGES ---
\usepackage{indentfirst} % 确保每一段的首行都会缩进
\setlength{\parindent}{2em} % 设置段首缩进为两个中文字符宽度
\setlength{\parskip}{0.3em} % 设置段落之间的间距

\begin{document}

\section*{导论}

\subsection*{习新思想创立的时代背景}

世界百年未有之大变局加速演进

中华民族伟大复兴进入关键时期

中国式现代化全面推进拓展

科学社会主义在21世纪的中国焕发新的蓬勃生机

中国共产党自我革命开辟新的境界

\subsection*{两个结合}

习新思想是马克思主义基本原理同中国具体实际相结合、同中华优秀传统文化相结合的重大成果。

\subsection*{习新思想的主要内容}

党的十九大、十九届六中全会提出的“十个明确”“十四个坚持”“十三个方面成就”

党的二十大提出的“六个必须坚持”

\subsection*{十个明确（部分）：}

明确中国特色社会主义最本质的特征是中国共产党领导，中国特色社会主义的最大优势是中国共产党领导，中国共产党是最高政治领导力量。

明确中国特色社会主义事业总体布局是经济建设、政治建设、文化建设、社会建设、生态文明建设五位一体。

\subsection*{六个必须坚持：}

必须坚持人民至上、必须坚持自信自立、必须坚持守正创新、必须坚持问题导向、必须坚持系统观念、必须坚持胸怀天下。

\subsection*{两个确立：}

党的十八届六中全会明确了习近平同志党中央的核心、全党的核心地位。

党的十九大把习近平新时代中国特色社会主义思想确立为党必须长期坚持的指导思想。

党的十九届六中全会指出，党确立习近平同志为党的核心、全党的核心地位，确立习新思想的指导地位。

\section*{第一章}

\subsection*{正确评价改革开放前后}

改革开放前后的历史不能互相否定

改革开放前的社会主义实践探索为改革开放后的社会主义实践探索积累了条件，后者是对前一个时期的坚持、改革、发展。

\subsection*{四个自信}

坚定道路自信、理论自信、制度自信、文化自信。

\subsection*{中国特色社会主义新时代}

中国特色社会主义进入新时代，是我国社会主要矛盾发生新变化的反应。

中国特色社会主义进入新时代，是党的主要任务发生新变化的反应。

中国特色社会主义进入新时代，是中国和世界关系发生新变化的反应。

\subsection*{社会主要矛盾}

党的十八大以来，经过科学分析，我们党及时做出我国社会主要矛盾已经转化为人民日益增长的美好生活需要和不平衡不充分发展之间的矛盾的重大战略判断。

新时代我国社会主要矛盾的变化，是在社会主义初级阶段中发生的变化，没有改变对我国社会主义所处历史阶段的判断。

我国仍处于并将长期处于社会主义初级阶段的基本国情没有变，我国是世界最大发展中国家的国际地位没有变。

\subsection*{四个全面}

全面建设社会主义现代化国家是战略目标，在“四个全面”中居于引领地位；

全面深化改革、全面依法治国、全面从严治党是三大战略举措。

\section*{第二章}

实现中华民族伟大复兴的中国梦，本质是国家富强、民族振兴、人民幸福。

\subsection*{中国式现代化的中国特色}

中国式现代化是人口规模巨大的现代化。

中国式现代化是全体人民共同富裕的现代化。

中国式现代化是物质文明和精神文明相协调的现代化。

中国式现代化是人与自然和谐共生的现代化。

中国式现代化是走和平发展道路的现代化。

\section*{第三章}

\subsection*{坚持党的全面领导}

中国共产党领导是中国特色社会主义最本质的特征

中国最大的国情就是中国共产党的领导

中国共产党领导是中国特色社会主义制度的最大优势

\subsection*{坚持党对一切工作的领导}

中国共产党是最高政治领导力量

党的领导是全面的、系统的、整体的

\section*{第四章}

人民性是马克思主义的本质属性和鲜明品格。

\subsection*{人民就是江山，江山就是人民}

中国共产党领导人民打江山、守江山，守的是人民的心。

人民是历史的创造者，是真正的英雄。

人民是创造历史的真正动力，是历史发展和社会进步的主体力量。

尊重人民历史地位，充分发挥人民主体作用。

民心是最大的政治，决定事业兴衰成败。

\subsection*{人民立场}

坚持人民立场，就要始终牢记党的初心和使命。

坚持人民立场，就要始终保持党同人民群众的血肉联系。

坚持人民立场，就要热爱人民、尊重人民、敬畏人民。

\subsection*{群众路线}

群众路线是我们党的生命线和根本工作路线

调查研究是获得真知灼见的源头活水，是贯彻群众路线的有效途径

\subsection*{推动全体人民共同富裕}

共同富裕是中国特色社会主义的本质要求，是中国式现代化的重要特征

实现共同富裕不仅是经济问题，而且是关系党的执政基础的重大政治问题

要从全局角度来把握共同富裕

扎实推进共同富裕，必须坚持正确的原则和科学的思路

推动全体人民共同富裕与促进人的全面发展是高度统一的

\section*{第五章}

\subsection*{改革开放}

改革开放是党和人民大踏步赶上时代的重要法宝

改革开放是坚持和发展中国特色社会主义的必由之路

全面深化改革开放是完成新时代目标任务的必然要求

中国要前进，就要全面深化改革开放

\subsection*{坚持全面深化改革开放的正确方向}

必须坚持和改善党的全面领导

必须坚持以人民为中心，促进社会公平正义、增进人民福祉

必须有利于进一步解放思想、进一步解放和发展社会生产力、进一步解放和增强社会活力

\subsection*{全面深化改革开放的总目标是}

完善和发展中国特色社会主义制度、推进国家治理体系和治理能力现代化

\subsection*{改革开放永无止境}

全面深化改革开放，是新时代坚持和发展中国特色社会主义的根本动力

改革开放永无止境是社会基本矛盾运动规律的深刻反映

改革开放永无止境是总结世界社会主义实践经验得出的重要结论

改革开放永无止境是推进党和人民事业发展的必然要求

\section*{第六章}

\subsection*{我国进入发展新阶段}

党的十八大以来，提出我国经济发展正处于增长速度换挡期、结构调整阵痛期、前期刺激政策消化期“三期叠加”阶段，进入了新常态。

新发展阶段是我国社会主义初级阶段历史进程中的一个重要阶段，是中国共产党带领人民迎来从站起来、富起来到强起来历史性跨越的新阶段，是全面建设社会主义现代化国家、向第二个百年奋斗目标进军的重要阶段。

\subsection*{贯彻新发展理念}

坚定不移贯彻创新、协调、绿色、开放、共享的新发展理念。

党的十八大以来，我们党适应发展新的形势，对发展理念和思路作出及时调整，提出了经济社会发展的许多重大理论和战略，其中新发展理念是最重要、最主要的。

\subsection*{高质量发展}

高质量发展，是能够很好满足人民日益增长的美好生活需要的发展

是体现新发展理念的发展

是创新成为第一动力、协调成为内生特点、绿色成为普遍形态、开放成为必由之路、共享成为根本目的的发展

\subsection*{社会主义基本经济制度}

公有制为主体、多种所有制经济共同发展；按劳分配为主体、多种分配方式并存；社会主义市场经济体制等

\subsection*{坚持两个毫不动摇}

毫不动摇巩固和发展公有制经济

毫不动摇鼓励、支持、引导非公有制经济发展

\subsection*{新发展格局}

构建以国内大循环为主体、国内国际双循环相互促进的新发展格局

大力推动构建新发展格局

着力推动实施扩大内需战略同深化供给侧结构性改革有机结合

着力发展实体经济

着力加快科技自立自强

着力推动产业链供应链优化升级

\section*{第七章}

经过长期努力，我国已建成世界上规模最大的教育体系

\subsection*{教育的意义}

教育是国之大计、党之大计

教育是国家经济社会发展的支撑力量

\subsection*{落实立德树人根本任务}

育人的根本在于立德，落实立德树人根本任务是建设教育强国的核心课题

着力解决好培养什么人的问题

着力解决要怎样培养人的问题

着力解决好为谁培养人的问题

增强自主创新能力，基础研究是科技创新的源头。

\section*{第八章}

制度优势是一个政党、一个国家的最大优势

\subsection*{人民民主是社会主义的生命}

民主是全人类共同价值，是人类政治文明发展的成果

民主是中国共产党和中国人民始终不渝坚持的重要理念

人民民主建立在社会主义经济基础之上，体现了社会主义国家的性质，反映了社会主义制度的本质要求，是一种新型的社会主义民主

人民民主是全面建设社会主义现代化强国的应有之义

\subsection*{中国特色社会主义政治制度是由根本政治制度、基本政治制度、重要政治制度等组成的制度体系}

根本：人民代表大会制度

基本：中国共产党领导的多党合作和政治协商制度、民族区域自治制度以及基层群众自治制度

\subsection*{坚定不移走中国特色社会主义政治发展道路}

必须坚持党的领导、人民当家作主、依法治国有机统一

要加强党中央集中统一领导，改进党的领导方式和执政方式，扩大人民有序政治参与，维护国家法制统一、尊严、权威，保障我国社会主义民主政治制度有序运行。

必须积极稳妥推进政治体制改革

必须始终保持政治定力

\subsection*{全过程人民民主}

全过程人民民主是过程民主和成果民主的统一

全过程人民民主是程序民主和实质民主的统一

全过程人民民主是直接民主和间接民主的统一

全过程人民民主是人民民主和国家意志的统一

全过程人民民主是全链条、全方位、全覆盖的民主

全过程人民民主是最广泛、最真实、最管用的民主

\subsection*{统一战线的工作范围包括}

民主党派成员

无党派人士

党外知识分子

少数民族人士

宗教界人士

非公有制经济人士

新的社会阶层人士

出国和归国留学人员

香港同胞、澳门同胞、台湾同胞及其在大陆的亲属

华侨、归侨及侨眷

其他需要联系和团结的人员

\section*{第九章}

全面依法治国的唯一正确道路是中国特色社会主义法治道路

\subsection*{中国特色社会主义法治道路的核心要义}

坚持党的领导

坚持中国特色社会主义制度

贯彻中国特色社会主义法治理论

\subsection*{坚持和拓展中国特色社会主义法治道路}

坚持中国共产党的领导

坚持以人民为中心

坚持法律面前人人平等

坚持依法治国和以德治国相结合

坚持从中国实际出发

\subsection*{统筹处理全面依法治国的重大关系}

正确处理政治和法治的关系

正确处理改革和法治的关系

正确处理依法治国和以德治国的关系

正确处理依法治国和依规治党的关系

\subsection*{全面推进依法治国的总抓手就是建设中国特色社会主义法治体系}

宪法是国家的基本法

全面推进科学立法、严格执法、公正司法、全民守法

\section*{第十章}

\subsection*{文化繁荣兴盛是实现中华民族伟大复兴的必然要求}

文化繁荣兴盛是实现中华民族伟大复兴的精神支撑

文化繁荣兴盛是建设社会主义现代化强国的应有之义

文化繁荣兴盛是满足人民日益增长的美好生活需要的内在要求

文化繁荣兴盛是在世界文化激荡中站稳脚跟的前提基础

文化自信是更基础、更广泛、更深厚的自信，是一个国家、一个民族发展中最基本、最深沉、最持久的力量

\subsection*{坚持马克思主义在意识形态领域指导地位的根本制度}

\subsection*{推动中华优秀传统文化创造性转化、创新性发展}

\section*{第十一章}

关注人本身的发展，是社会主义和资本主义的根本区别之一

\subsection*{坚持在发展中增进民生福祉}

发展是解决民生问题的“总钥匙”，民生是发展的“指南针”

正确把握民生和发展的关系，是保障和改善民生的重要前提

坚守底线、突出重点、完善制度、引导预期，是保障和改善民生的工作思路

解决人民群众最关心最直接最现实的利益问题，是保障和改善民生的重中之重

坚持人人尽责、人人享有，让所有劳动者在推动发展中分享发展成果，是保障和改善民生的重要原则

\subsection*{就业是最基本的民生}

党的十八届三中全会提出创新社会治理体制，从社会管理到社会治理，意味着社会建设理念、体制和方式的重大变革，体现了我们党对社会运行和社会治理规律认识的不断深化。

\section*{第十二章}

\subsection*{习近平生态文明思想的主要内容集中体现为“十个坚持”}

坚持党对生态文明建设的全面领导

坚持生态兴则文明兴

坚持人与自然和谐共生

坚持绿水青山就是金山银山

坚持良好生态环境是最普惠的民生福祉

坚持绿色发展是发展观的深刻革命

坚持统筹山水林田湖草沙系统治理

坚持用最严格制度最严法治保护生态环境

坚持把建设美丽中国转化为全体人民自觉行动

坚持共谋全球生态文明建设之路

\subsection*{中华民族向来尊重自然、热爱自然}

而不是畏惧自然，可以说敬畏自然

\subsection*{绿水青山就是金山银山}

生态环境保护和经济发展二者是辩证统一、相辅相成的关系

绿水青山既是自然财富、生态财富，又是社会财富、经济财富

绿水金山可带来金山银山，金山银山买不到绿水青山

绿水青山还是人民群众健康的重要保障，是人民群众的共有财富

坚持生态优先、绿色发展，把生态环境优势转化为生态农业、生态工业、生态旅游等生态经济优势，绿水青山就变成了金山银山

处理好绿水青山和金山银山的关系，关键在人，关键在思路

2030前碳达峰，2060前碳中和

\section*{第十三章}

\subsection*{总体国家安全观是新时代国家安全工作的基本遵循}

习近平在中央国家安全委员会第一次会议上首次提出总体国家安全观，强调必须坚持总体国家安全观，走出一条中国特色国家安全道路。

\subsection*{总体国家安全观以人民安全为宗旨}

以政治安全为根本，以经济安全为基础，以军事、科技、文化、社会安全为保障，以促进国际安全为依托

\subsection*{总体国家安全观的关键是“总体”}

总体国家安全观，运用总体战略思维和宽广世界眼光把握国家安全，既回答了如何解决好大国发展进程中面临的共性安全问题，同时又回答了如何解决好中华民族伟大复兴关键时期面临的共性安全问题

\subsection*{总体国家安全观十个坚持}

坚持党对国家安全工作的绝对领导

坚持中国特色国家安全道路

坚持以人民安全为宗旨

坚持统筹发展和安全

坚持把政治安全放在首要位置

坚持统筹推进各领域安全

坚持把防范化解国家安全风险摆在突出位置

坚持推进国际共同安全

坚持推进国家安全体系和能力现代化

坚持加强国家安全干部队伍建设

\subsection*{统筹发展和安全}

统筹发展和安全，是党和国家的一项基础性工作，是我们党治国理政的一个重大原则

发展和安全是两件大事

发展解决的是动力问题，是推动国家和民族赓续绵延的根本支撑；安全解决的是保障问题，是确保国家和民族行稳致远的坚强柱石

发展具有基础性、根本性，是解决安全问题的总钥匙，发展就是最大的安全

安全是发展的条件和保障

必须同步推进发展和安全。要坚持发展和安全并重，把国家安全同经济社会发展一起谋划、一起部署

以新安全格局保障新发展格局。牢牢守住安全发展这条底线是构建新发展格局的重要前提和保障

主动塑造于我有利的外部安全环境，推动发展和安全深度融合，实现高质量发展和高水平安全的良性互动

发展是安全的基础，安全是发展的条件——习近平

\subsection*{把维护政治安全放在首要位置}

政治安全决定和影响着国家经济安全、军事安全、社会安全等各个领域的安全

维护政权安全，就是要毫不动摇坚持和巩固党的领导和长期执政地位

维护制度安全，就是要毫不动摇坚持和完善中国特色社会主义制度

维护意识形态安全，就是要毫不动摇坚持和巩固马克思主义在意识形态领域的指导地位，不断巩固全党全国人民团结奋斗的共同思想基础

政治安全与人民安全、国家利益至上是有机统一的

\subsection*{重点领域国家安全}

国土安全

经济安全

社会安全。社会安全与人民群众切身利益关系最密切，是人民群众安全的晴雨表，是社会安定的风向标

网络、人工智能、数据安全

生物安全和公共卫生安全

外部安全

\subsection*{提高防范化解重大风险能力}

坚持底线思维和极限思维

下好先手棋、打好主动仗，力争把风险化解在源头

运用制度威力应对风险挑战的冲击

\section*{第十四章}

党的十九大明确提出，党在新时代的强军目标是建设一支听党指挥，能打胜仗，作风优良的人民军队，把人民军队建设成为世界一流军队

\subsection*{新时代人民军队使命任务}

为巩固中国共产党领导和我国社会主义提供战略支撑

为捍卫国家主权、统一和领土完整提供战略支撑

为维护我国海外利益提供战略支撑

为促进世界和平和发展提供战略支撑

\subsection*{建设强大军队}

建设强大的人民军队是我们党的不懈追求

习近平指出，建设强大军队是接续奋斗的伟大事业

\subsection*{强军目标的科学内涵}

听党指挥是灵魂，决定军队建设的政治方向。坚持党对人民军队的绝对领导是人民军队永远不变的军魂

能打胜仗是核心，反映军队的根本职能和军队建设的根本指向

作风优良是保障，关系军队的性质、宗旨、本色

\subsection*{坚持党对人民军队的绝对领导}

坚持党对人民军队的绝对领导，是马克思主义建党建军的一条基本原则

党对人民军队的绝对领导，是建军之本、强军之魂

坚持党对人民军队的绝对领导必须有一整套制度作保障

军委主席负责制是坚持党对人民军队绝对领导的根本制度和根本实现形式，在党领导军队的一整套制度体系中处于最高层次、居于统领地位

坚持政治建军、改革强军、科技强军、人才强军、依法治军

\end{document}
